\section{An exact split cubic in closed H sewing}\label{sec:H}

Let \(M_H\) be the six-component H plumbing: join two zero-framed centres by one Hopf edge; at each centre attach the two meridian leaves of framings \(+2,-2\). All nonzero linking magnitudes are one, and no other components link. We use positive Hopf convention for the spine, equivalent to a signed exchange of the induced peripheral axes. The final scalar here is independent of this sign. Denote by \(\rho_H\) the quaternionic character described below.

\begin{theorem}[Localized closed logarithm]\label{thm:H}
At \(\rho_H\), the harmonic Coulomb-reduced Chern--Simons action is exactly a nonzero split real cubic \(cXYZ\), and \(\CS(\rho_H)=1/2\pmod1\). In arbitrarily small fixed neighborhoods of this chart the original unmarked closed level sum has the representation
\begin{equation}\label{eq:localH}
 Z_r(M_H)=(-1)^r r^{3/2}\!
   \int_{\R^3} e^{2\pi\ii r cXYZ}\,\varrho(X,Y,Z)\,dX\,dY\,dZ
     +O(\sqrt r),
 \qquad \varrho(0)=\frac{\sqrt2\,\operatorname{Vol}(M_H)}{2\pi^3}.
\end{equation}
Here \(\varrho\) is a smooth local density derived from the surgery expression with fixed smooth cutoffs. One may take \(|c|=\operatorname{Vol}(M_H)/\pi^2\) in orthonormal parallel coordinates. Consequently the entire leading term of the closed scalar is \((-1)^r\sqrt{2r}\log r/\pi\), or \((-1)^r r^2\log r/\pi^2\) for \(\tau_r\). Formula \eqref{eq:localH} does not identify the \(O(\sqrt r)\) part.
\end{theorem}

\subsection{Classical chart and the action value}

Each half of the H graph has generating sections \(x,y,z\), fibre \(h\), and relations \(x^2 h=1,\ y^{-2}h=1,\ z=(xy)^{-1}\), after choices of cone meridian signs. There is a marked flat model
\[
 (\R^2\times[-1,1])/\langle A,B\rangle,\quad
 A(X_0,Y_0,t)=(X_0+1,-Y_0,-t),\quad B(X_0,Y_0,t)=(X_0,Y_0+2,t).
\]
The circle fibration in the \(X_0\) direction has disk base with two order-two cone points. The two cone lifts can be chosen as \(x=A,\ y=A^{-1}B\) with \(h=A^{-2}\); indeed their squares give exactly the marked filling slopes, and the product section is \(z=B^{-1}\). The two lifts can be prescribed independently because a section change over the pair of pants has two independent integral circle degrees. This checks the actual peripheral marking, not just its abstract group. On the boundary the exchanged directions have equal periods and product collars, so the marked exchange is a flat isometry. This is the flat torus semi-bundle case~\cite{SWW}. Writing \(C\) for the other glide, van Kampen gives \(AC^2A^{-1}=C^{-2},\ CA^2C^{-1}=A^{-2}\). Sending \(A,C\) to orthogonal imaginary unit quaternions \(\mathbf i,\mathbf j\) gives a representation. Their adjoint half-turns are exactly the tangent linear \(V_4\) representation.

Here is why its cubic is exact. For a flat connection with \(*[\mathcal H^1,\mathcal H^1]\subset\mathcal H^1\), split the Sobolev Coulomb slice as \(a+b\) with \(a\) harmonic and \(b\) coexact. The projected transverse stationarity equation is
\[
 P_{\perp}*_3\big(d_A b+\tfrac12[a+b,a+b]\big)=0 .
\]
Its coexact linearization \(*d_A\) is an isomorphism from order \(s\) to order \(s-1\), \(s>3/2\): its square is \(d_A^*d_A\) on this summand, with no harmonic kernel. The implicit function theorem gives a unique small \(b(a)\); the harmonic bracket condition makes \(b=0\) an exact solution. On a flat compact manifold with flat unitary coefficients, Bochner's identity makes twisted harmonic forms parallel, hence the bracket condition holds. Harmonic one-forms are equivariant linear maps between the tangent and adjoint representations. Both \(V_4\) representations split into the same three distinct real nontrivial characters, so the harmonic form is \(X\mathbf i\,dx_1+Y\mathbf j\,dx_2+Z\mathbf k\,dx_3\). The cubic trace has \(\operatorname{Tr}(\mathbf i\mathbf j\mathbf k)=-2\), giving \(|c|=\operatorname{Vol}/\pi^2\); there is no quadratic term. This representation has no infinitesimal stabilizer. The exact critical set is the three axes. The argument identifies a reduced action, not a measure on all connections.

For the action value, rotate the second generator from \(\mathbf j\) to \(\mathbf i\) with both squares fixed at \(-1\), an actual flat path. The endpoint eigenline sends both plus-cap meridians to \(1/4\in\R/\Z\). The surgery intersection matrix, on the two centres followed by the two pairs of leaves, is
\[
 \begin{pmatrix}0&1&1&1&0&0\\1&0&0&0&1&1\\1&0&2&0&0&0\\
 1&0&0&-2&0&0\\0&1&0&0&2&0\\0&1&0&0&0&-2\end{pmatrix}.
\]
Elimination gives \(H_1=(\Z/4)^2\) generated by the two plus meridians, with linking matrix \(\bigl(\begin{smallmatrix}1/2&1/4\\1/4&1/2\end{smallmatrix}\bigr)\) up to overall orientation. The corresponding linking dual has square \(1/2\) modulo one. Since \(c_2(L\oplus L^{-1})=-c_1(L)^2\), the fundamental-trace transgression of this eigenline gives \(\CS=1/2\pmod1\), unaffected by an overall sign; see~\cite{Auckly}. Constancy along the flat path gives the irreducible action value.

\subsection{Finite localization and density normalization}

From \eqref{eq:rec},
\[
 |C_2(r-n)|^2=\tfrac4r{\bf1}_{n\ {\rm odd}}\cos^2(\pi n/(2r)).
\]
The exact H scalar is thus
\[
 Z_r=(-1)^r(2/r)^{3/2}\sum_{n,m\ {\rm odd}<r}
      \cot(\pi n/(2r))\cot(\pi m/(2r))\sin(\pi nm/r).
\]
For reference the finite sine transform
\[
 \tfrac2r\sum_{n\ {\rm odd}<r}\cot(\pi n/(2r))\sin(\pi nm/r)=1,\qquad 1\le m<r
\]
follows by doubling to the full odd grid modulo \(2r\) and using
\(\sin(mq)\cot(q/2)=1+\cos(mq)+2\sum_{j=1}^{m-1}\cos(jq)\). In particular the exact scalar also equals \((-1)^r\sqrt{2/r}\sum_{m\ {\rm odd}<r}\cot(\pi m/(2r))\). The orbital expression is stronger for localization. Let \(x=n/r,y=m/r,\ h=2/r\), \(a(x)=x\cot(\pi x/2)\), \(a(0)=2/\pi\). Since the relative dot product \(t\) of two Haar-transported adjoint axes is uniform on \([-1,1]\), the edge sine becomes an exact orbital integral, giving
\begin{equation}\label{eq:Horbit}
 Z_r=(-1)^r\frac{\pi\sqrt2}{4}r^{3/2}
 h^2\!\sum_{x,y\ {\rm positive\ odd\ mesh}<1}\!
 a(x)a(y)\int_{-1}^{1}e^{\pi\ii r xyt}\,dt.
\end{equation}

Choose even smooth cutoffs \(\chi(x)\chi(y)\eta(t)\), all one near zero and supported in arbitrarily small fixed intervals. The part omitted from \eqref{eq:Horbit} is \(O(\sqrt r)\). For the exterior charge part, the exact sine transform is the odd \(2r\)-periodic square wave, with zero at its breaks. Multiplication by an even smooth period-two weight \(f(n/r)\) with \(f(0)=0\) convolves it with rapidly decreasing Fourier coefficients. Away from a break the constant terms sum to zero, so the new transform at \(m\) is \(O_N((1+\min(m,r-m))^{-N})\). The remaining cotangent, bounded by \(O(r/m)\) at one end and \(O((r-m)/r)\) at the other, gives the claimed error. This works whenever one cutoff removes its near-zero charges. After both colour cutoffs, signed extension of the two odd grids and even \(t\)-integration converts their positive quadrant to one quarter of the full signed grid. Two-dimensional Poisson summation has phases \(\pi r(txy-u x-vy)\), \(u,v\in\Z\). All nonzero aliases are uniformly nonstationary since the colour support has \(|x|,|y|<d<1\) and \(|t|\le1\); repeated integration by parts sums their errors to \(O_N(r^{-N})\) before the prefactor. The zero alias with \(t\) away from zero is \(O(r^{-1})\) by integrating in \(x\) and using the Schwartz transform in \(ty\). This also removes the angular complement at \(O(\sqrt r)\). The remaining integral is
\[
 (-1)^r\,\frac{\pi\sqrt2}{16}r^{3/2}
 \int_{\R^3} a(x)\chi(x)a(y)\chi(y)\eta(t)e^{\pi\ii r xyt}\,dx\,dy\,dt .
\]

To calibrate its variables, let the spatial glide translations be \((a_0,0,0)\) for \(A\) and \((0,b_0,d_0)\) for \(C\) after choosing a Euclidean origin on the first axis. The squares of \(A,C,AC\) generate the translation lattice: they span the respective independent directions, and the quotient generated by the two resulting involutions with involutive product has order four, the already known linear holonomy order. Thus \(\operatorname{Vol}=2|a_0b_0d_0|\). On the parallel chart the fixed straight representative paths yield holonomies
\[
 U=\mathbf i e^{a_0X\mathbf i},\qquad V=\mathbf j e^{b_0Y\mathbf j+d_0Z\mathbf k},
\]
up to a simultaneous choice of connection signs. Their regular fibres are \(U^{-2},V^{-2}\). The Wilson eigenvalue \(2\cos(\pi a/r)\) identifies the two signed conjugacy-defect lifts with \(x,y\). The third coordinate is the relative dot product \(t\) of the normalized axes of \(U,V\), calibrated by the exact orbital integral. Direct quaternion multiplication gives linear terms \(2a_0X/\pi,\,2b_0Y/\pi,\,d_0Z\) up to signs, and each of \(x,y,t\) vanishes exactly on its corresponding coordinate plane locally. Hence \(2\pi cXYZ=\pi xyt\,g(x,y,t)\) with \(g(0)=\pm1\) by the volume identity. Absorb \(g\) in the local \(t\) coordinate; the Jacobian at zero has absolute value one for this last change. The pushed density at zero is
\(\frac{\pi\sqrt2}{16}(2/\pi)^2\, 4|a_0b_0d_0|/\pi^2\), as in \eqref{eq:localH}. The phase \(r\) versus \(k=r-2\) at \(\CS=1/2\) is the same; the two actions in the local integral differ only by a smooth density multiplier with value one at zero.

Finally for any smooth compactly supported \(g_0\), the integral of \(e^{\ii\lambda cxyz}g_0\) has leading coefficient \(4\pi g_0(0)\log\lambda/(\lambda|c|)\). To see the remainder, subtract an even product cutoff times \(g_0(0)\); its error is a sum of coordinate multiples and has integral \(O(\lambda^{-1})\) by the Fourier bound \(C_N(1+\lambda|xy|)^{-N}\) in the third variable. The product term on a positive square yields \((\lambda|c|)^{-1}\int_0^{\lambda|c|d^2}\widehat\eta(s)\log(\lambda|c|d^2/s)\,ds\), and the four quadrants give the result since \(\int_0^\infty\widehat\eta(s)\,ds=\pi\eta(0)\). Substituting \(\lambda=2\pi r\) proves Theorem~\ref{thm:H}. The angular variable in this proof is an exact orbital variable, not a claim of a finite operator projecting onto arbitrary relative angles.
